%% -*- coding: utf-8 -*-
%% CA2 Progress Report
%% Higher Degree by Research -- Candidature Assessment 2 (CA2)
%% Standalone document; compiles independently of the STOP paper.

\documentclass[11pt]{article}

\usepackage[margin=1in]{geometry}
\usepackage{graphicx}
\usepackage{booktabs}
\usepackage{amsmath}
\usepackage{amssymb}
\usepackage{xcolor}
\usepackage{enumitem}
\usepackage[colorlinks=true, linkcolor=black, citecolor=black, urlcolor=blue]{hyperref}

\graphicspath{{figures/}}

% Placeholder macro for details to be filled in by the candidate.
\newcommand{\TODO}[1]{\textcolor{red}{\textbf{[TODO: #1]}}}

\title{\textbf{Candidature Assessment 2 (CA2) --- Progress Report}\\[0.4em]
\large Advances in Uncertainty Quantification for Machine Learning through\\
Diversity-Aware Learning and Efficient Inference}
\author{Candidate: Haowei Huang \\ Principal Supervisor: Dr. Junyu Xuan \quad }

% \author{Candidate: Haowei Huang \\ Supervisor: Dr. Junyu Xuan}

% \institute[]{Australia Artificial Intelligence Institute (AAII) \par
% Faculty of Engineering and Information Technology \par
% University of Technology Sydney, Australia \par} 

\begin{document}
\maketitle

\begin{abstract}
\noindent
This report documents progress on my doctoral project since CA1. The thesis advances
\emph{uncertainty quantification (UQ) for machine learning} along two complementary directions:
\emph{diversity-aware learning} --- controlling how ensemble members disagree so that
uncertainty is well calibrated --- and \emph{efficient inference} --- reducing the cost of
building and deploying ensembles. Within this scope, the report focuses on two contributions.
The first, STOP, is a substantially completed method that uses the edge-of-stability transition
in training dynamics to switch off diversity regularization once it is no longer needed, cutting
training cost while preserving accuracy and calibration. The second carries the same
diversity lens to large language models (LLMs): instead of \emph{controlling} member diversity
during training, it \emph{reads out} the structural diversity of an ensemble of Chain-of-Thought
reasoning chains and turns it into a black-box uncertainty estimate, adding no model inference
beyond the sampled chains. For each of these two completed contributions I state the method, the
key results, and how they bear on the central research question. A planned third contribution on
uncertainty for LLM reasoning, together with a detailed plan towards thesis completion --- timeline,
risks, and chapter structure --- closes the report. It is intended as the basis for discussion with my supervisory
team and the assessment panel.
\end{abstract}

% =====================================================================
\section{Research Aims and Questions}
\label{sec:aims}

\paragraph{Thesis scope.}
The thesis advances uncertainty quantification (UQ) for machine learning through two
complementary directions. \emph{Diversity-aware learning} concerns how ensemble members are
encouraged to disagree so that predictive uncertainty is well calibrated. \emph{Efficient
inference} concerns reducing the cost of building and using ensembles, which otherwise grows with
ensemble size and model scale. Deep ensembles~\cite{lakshminarayanan2017simplescalablepredictiveuncertainty}
are the vehicle throughout: they are among the most reliable UQ methods, precisely because members
explore distinct regions of the loss landscape~\cite{fort2020deepensembleslosslandscape}, but this
reliability is bought with substantial compute. This report covers two contributions under that
scope --- STOP (Section~\ref{sec:work1}) and a diversity-based uncertainty method for
LLMs (Section~\ref{sec:work2}).

\paragraph{Central research question.}
\emph{Can diversity be turned into well-calibrated, low-cost uncertainty --- during training, by
using ensemble dynamics to decide when diversity regularization is still needed; and at inference
for large language models, by reading the structural diversity of their reasoning chains?}

\paragraph{Research questions.}
This central question decomposes into two research questions, one per contribution in this report:
\begin{description}[leftmargin=2.4em, style=nextline, itemsep=2pt]
  \item[RQ1.] Does deep-ensemble training separate into an exploration and an exploitation phase
  around the edge of stability, and if so, can a dynamics-triggered rule switch off diversity
  regularization at that transition --- cutting training cost while preserving accuracy and
  calibration? \emph{(Addressed by Work~1, Section~\ref{sec:work1}.)}
  \item[RQ2.] Does the structural diversity of an LLM's Chain-of-Thought reasoning chains --- how
  central or outlying each chain is within a sampled ensemble --- provide a stronger black-box
  uncertainty signal than final-answer agreement, without any additional model inference?
  \emph{(Addressed by Work~2, Section~\ref{sec:work2}.)}
\end{description}

\paragraph{Objectives.}
\begin{description}[leftmargin=2.4em, style=nextline, itemsep=2pt]
  \item[O1.] Characterize deep ensemble training as a two-phase (exploration / exploitation)
  process separated by the edge of stability (EoS), and establish that useful diversity is
  formed in the exploration phase. \emph{(Diversity-aware learning.)}
  \item[O2.] Design and validate a practical, dynamics-triggered training strategy that reduces
  the cost of diversity regularization while preserving accuracy and calibration.
  \emph{(Efficient inference.)}
  \item[O3.] Carry the diversity lens to large language models: quantify uncertainty by measuring
  the structural diversity of an ensemble of Chain-of-Thought reasoning chains, under a black-box,
  model-free constraint. \emph{(Diversity-aware learning at inference.)}
\end{description}

Objectives O1--O2 are addressed by the STOP work (Section~\ref{sec:work1}) and O3 by the LLM
contribution (Section~\ref{sec:work2}), both now complete. A planned third contribution (Work~3),
extending the LLM-uncertainty line of O3, and the plan towards thesis completion are set out in
Section~\ref{sec:plan}.


% =====================================================================
\section{Work 1 --- STOP: Stability-Triggered Optimization for Deep Ensembles}
\label{sec:work1}

To make members disagree in useful ways, diversity regularizers such as
Repul~\cite{dangelo2021repulsive} and DICE~\cite{rame2021dice} add overhead that grows with
ensemble size. Table~\ref{tab:cost} illustrates this on a representative UCI dataset --- diversity
regularization inflates training time by up to $10\times$ at $M=20$ relative to a plain deep
ensemble (DE). The question driving STOP is whether this cost is \emph{necessary throughout
training}, or whether the training dynamics themselves reveal when diversity has already been
established and regularization can be safely stopped.

\begin{table}[ht]
  \centering
  \caption{Training time (seconds) on Concrete for increasing ensemble size $M$, showing the
  cost of diversity regularization relative to a plain deep ensemble (DE). Multipliers are
  relative to DE at the same $M$.}
  \label{tab:cost}
  \begin{tabular}{lccc}
    \toprule
    & \multicolumn{3}{c}{Ensemble size $M$} \\
    \cmidrule(lr){2-4}
    Method & $M=5$ & $M=10$ & $M=20$ \\
    \midrule
    DE    &  23.3\,s          &  51.0\,s          &   93.7\,s \\
    Repul &  51.6\,s\ (2.2$\times$) & 179.0\,s\ (3.5$\times$) & 384.7\,s\ (4.1$\times$) \\
    DICE  &  86.0\,s\ (3.7$\times$) & 269.2\,s\ (5.3$\times$) & 937.0\,s\ (10.0$\times$) \\
    \bottomrule
  \end{tabular}
\end{table}

\subsection{Research question and idea}
Single-model optimization with gradient descent is known to enter an \emph{edge of stability}
(EoS) regime~\cite{cohen2021gradient}, where the largest Hessian eigenvalue (sharpness)
$\lambda_{\max}$ hovers near the stability threshold $2/\eta$ for learning rate $\eta$. Work~1
extends this notion to deep ensembles. The hypothesis is that ensemble training decomposes into two phases:
\begin{itemize}[itemsep=2pt]
  \item an \textbf{exploration phase} before EoS, where members are dynamically unstable, rapidly
  diversify, and settle into distinct basins --- and where diversity regularization matters;
  \item an \textbf{exploitation phase} after EoS, where each member refines its solution within an
  already-selected basin --- and where continued regularization yields diminishing returns at
  full cost.
\end{itemize}
If true, one should regularize \emph{only} until EoS and stop thereafter. This is STOP:
Stability-Triggered Optimization. Per member $k$, sharpness is monitored (after a short warm-up)
and diversity regularization is disabled once member $k$'s sharpness first exceeds a scaled
threshold; the ensemble is considered past EoS when all members have triggered.

\subsection{Evidence that the two-phase picture holds}
Figure~\ref{fig:trajectory} gives the motivating intuition in a controlled 2D setting: without
regularization two members collapse into the same basin; with always-on regularization they
separate; with STOP they \emph{still} separate despite regularization being switched off early ---
consistent with basin selection becoming effectively irreversible past EoS.

\begin{figure}[ht]
  \centering
  \includegraphics[width=0.92\linewidth]{learning_trajectory_2d_comparison.png}
  \caption{Optimization trajectories of ensemble members. Left: no regularization (members
  collapse). Middle: always-on diversity regularization (members separate). Right: STOP ---
  regularization active only before EoS, then disabled; members still reach distinct basins.}
  \label{fig:trajectory}
\end{figure}

Figure~\ref{fig:diversity} provides the direct empirical test on real data (pumadyn, 5 members),
added in response to reviewer feedback. Sharpness rises and saturates near $2/\eta$, and the
diversity signals (function-space distance, Jensen gap) plateau around the per-member EoS
triggers --- i.e., diversity is essentially formed by the time STOP switches regularization off.

\begin{figure}[ht]
  \centering
  \includegraphics[width=0.92\linewidth]{pumadyn_dice_repul_pair_diversity.png}
  \caption{Sharpness and diversity trajectories on pumadyn (5 members): DICE (left), Repul
  (right). Rows: sharpness $\lambda_{\max}$; function-space distance; RMSE; Jensen gap. Vertical
  lines mark per-member EoS triggers. Diversity plateaus around the triggers, supporting the
  two-phase account.}
  \label{fig:diversity}
\end{figure}

\subsection{Key results}
\paragraph{Efficiency (the headline result).}
On CIFAR-10 with ResNet-20~\cite{he2016deep} (Table~\ref{tab:cifar_time}), STOP cuts total training time by
57--66\% for DICE and 59--73\% for Repul, with savings \emph{growing} with ensemble size ---
mirroring the quadratic scaling of regularization cost. A component-level breakdown (UCI) shows
sharpness monitoring costs only 2--5\% of total training time, so EoS detection is cheap relative
to the regularization it removes.

\begin{table}[ht]
\centering
\caption{Total training time (hours, $\downarrow$) on CIFAR-10 (ResNet-20). \textbf{Bold} = faster.}
\label{tab:cifar_time}
\small
\begin{tabular}{lrrrrrr}
\toprule
Method & \multicolumn{2}{c}{$M=5$} & \multicolumn{2}{c}{$M=10$} & \multicolumn{2}{c}{$M=20$} \\
\cmidrule(lr){2-3} \cmidrule(lr){4-5} \cmidrule(lr){6-7}
 & Base & +STOP & Base & +STOP & Base & +STOP \\
\midrule
DICE  & 21.7 & \textbf{9.3} & 32.9 & \textbf{12.3} & 129.4 & \textbf{44.1} \\
Repul & 4.1  & \textbf{1.7} & 8.4  & \textbf{3.5}  & 25.2  & \textbf{6.8}  \\
\bottomrule
\end{tabular}
\end{table}

\paragraph{Predictive quality is preserved.}
On UCI regression (Table~\ref{tab:uci}), STOP is no worse than always-on regularization on MSE
and generally matches or improves calibration (QCE). On CIFAR-10, clean accuracy stays within
$\sim$0.4\% of the always-on baselines and STOP frequently \emph{improves} robustness under
input corruption --- i.e., the large efficiency gains do not come at the expense of accuracy or
uncertainty quality.

\begin{table}[ht]
\centering
\caption{MSE ($\downarrow$) and QCE ($\downarrow$) on UCI datasets ($M=5$). \textbf{Bold} = STOP
is no worse than the corresponding always-on method.}
\label{tab:uci}
\small
\setlength{\tabcolsep}{4pt}
\begin{tabular}{l rrrrr rrrrr}
\toprule
& \multicolumn{5}{c}{MSE ($\downarrow$)} & \multicolumn{5}{c}{QCE ($\downarrow$)} \\
\cmidrule(lr){2-6} \cmidrule(lr){7-11}
Dataset & DE & DICE & DICE+STOP & Repul & Repul+STOP & DE & DICE & DICE+STOP & Repul & Repul+STOP \\
\midrule
Pumadyn  & 0.12  & 0.07  & \textbf{0.07}  & 0.11  & \textbf{0.09}  & 0.15 & 0.12 & \textbf{0.08} & \textbf{0.04} & \textbf{0.04} \\
Concrete & 27.00 & 22.06 & \textbf{21.84} & 26.57 & \textbf{21.84} & 0.25 & 0.21 & \textbf{0.14} & 0.09 & \textbf{0.08} \\
Airfoil  & 5.60  & 5.43  & \textbf{5.01}  & 4.85  & \textbf{4.81}  & 0.20 & 0.19 & \textbf{0.18} & \textbf{0.14} & \textbf{0.14} \\
Slump    & 39.00 & 37.36 & \textbf{36.14} & 37.93 & \textbf{37.85} & 0.25 & 0.22 & 0.22 & \textbf{0.16} & \textbf{0.16} \\
\bottomrule
\end{tabular}
\end{table}

\paragraph{The trigger, not merely early stopping, drives the gain.}
Baseline comparisons over fixed-stop schedules (30/50/70/90\%) and a matched-epoch stop with no
sharpness computation show that \emph{EoS-aware} triggering --- not simply regularizing for fewer
epochs --- is what preserves accuracy and calibration at the reduced cost, confirming that the
edge-of-stability signal, and not an arbitrary early cutoff, is doing the work.

\subsection{Status and reviewer response}
The STOP paper has been written and submitted (\TODO{venue / date}). Reviewer feedback centred on
(i) direct evidence for the two-phase claim, (ii) stronger baselines, (iii) sensitivity of the
trigger, and (iv) honest runtime accounting including monitoring overhead. In response I completed
four experiment sets, building on the ensemble UQ evaluation pipeline of~\cite{seligmann2024beyond}:
\textbf{(A)} sharpness/diversity trajectory plots (Fig.~\ref{fig:diversity});
\textbf{(B)} fixed-stop and matched-stop baselines; \textbf{(C)} sensitivity ablations over
threshold and cadence; \textbf{(D)} a wall-clock runtime breakdown isolating sharpness-monitoring
cost. These directly address the main criticisms and are being folded into a revised manuscript.

% =====================================================================
\section{Work 2 --- Reading Out Reasoning Diversity for Uncertainty in Large Language Models}
\label{sec:work2}

Work~2 carries the \emph{diversity} lens of the thesis from training-time \emph{control} to
inference-time \emph{measurement}. Where Work~1 asks how to make ensemble members disagree in useful
ways --- and when to stop paying for it --- Work~2 asks the dual question for large language models
(LLMs): given an ensemble of reasoning chains a single LLM produces for one question, how much do
they \emph{structurally} agree, and can that agreement be turned into a calibrated confidence score?
The setting is black-box --- only input/output access, no logits or hidden states --- so the
ensemble is formed by sampling multiple Chain-of-Thought (CoT) responses rather than by training
$M$ separate models. Efficiency therefore enters as a \emph{constraint the method respects} (no LLM
calls beyond the sampled ensemble; only a small sentence encoder), not as the object of study.

\subsection{Motivation and idea}
Existing black-box UQ for LLMs scores confidence from \emph{final-answer} agreement across sampled
responses --- self-consistency~\cite{wang2023selfconsistency}, semantic
entropy~\cite{kuhn2023semantic,farquhar2024detecting}, and verbalized / self-probing
confidence~\cite{xiong2024llms}. This collapses each response to its answer and discards the
reasoning that produced it. Two chains that reach the \emph{same} answer through \emph{different}
reasoning represent genuine epistemic uncertainty that answer-level methods cannot see
(Figure~\ref{fig:intro_motivation}); conversely,
chains that converge on a shared reasoning \emph{structure} are more trustworthy than a bare answer
count suggests. CoT models expose this intermediate structure explicitly~\cite{wei2022chain}, and
recent work has begun to exploit reasoning steps for UQ~\cite{zhang2025cotuq}. Work~2 treats the
ensemble of reasoning chains as a \emph{sample of curves} and asks which chains are central and
which are outlying --- exactly the question functional-data \emph{band
depth}~\cite{lopezpintado2009banddepth} was designed to answer. Centrality of a reasoning chain
within the ensemble becomes the confidence signal, making Work~2 a direct, inference-time
counterpart to the diversity-aware learning of Work~1.

\begin{figure}[ht]
  \centering
  \includegraphics[width=0.98\linewidth]{intro_motivation}
  \caption{Answer-level versus reasoning-level uncertainty. Two ensembles of three CoT chains split
  their final answers identically (two answer $B$, one answers $A$), so answer-level agreement
  reports the same confidence and cannot tell them apart. What differs is \emph{where} the chains
  diverge --- a late split (near-agreement, left) versus an early split (genuine uncertainty,
  right). Band depth scores exactly this reasoning structure, which answer-level agreement discards.}
  \label{fig:intro_motivation}
\end{figure}

\subsection{Method}
Each sampled response is parsed into an ordered sequence of reasoning steps, and every step is
embedded with a small sentence encoder~\cite{reimers2019sentencebert}. A response's confidence is
its \emph{centrality} in the ensemble of chains: central (high-depth) chains lie ``inside'' the
bundle the others trace out and are treated as confident, while outlying chains signal uncertainty.
I develop a family of estimators of increasing structural fidelity
(Figure~\ref{fig:methods_overview}):
\begin{itemize}[itemsep=2pt]
  \item \textbf{pBD (path band depth).} Treats each chain as an ordered curve, aligns chains of
  different lengths by dynamic time warping, and scores a chain by the mean similarity of its steps
  to the ensemble at aligned positions. Order-sensitive.
  \item \textbf{vBD (vertex band depth).} Merges semantically equivalent steps across chains into a
  reasoning graph and scores each chain by the average band depth of the graph vertices it visits.
  Order-invariant, but path structure is lost.
  \item \textbf{gBD (graph band depth).} Restores path-level depth but measures it against
  \emph{random walks on the reasoning graph} rather than the handful of sampled chains, thereby
  characterizing the full reasoning space the graph defines --- not only the chains the model
  happened to emit. This is the strongest variant.
\end{itemize}
All three are model-free: they need only the sampled ensemble and a sentence encoder, and add no
LLM inference beyond generating the chains.

\begin{figure}[ht]
  \centering
  \includegraphics[width=0.98\linewidth]{methods_overview}
  \caption{The three band-depth estimators, viewed as bands of curves. Each curve is one sampled
  CoT chain and each node a reasoning step; a chain is \emph{deep} when it lies inside the band
  formed by the others. \textbf{pBD} keeps the chains as position-aligned curves (order-sensitive);
  \textbf{vBD} collapses semantically equivalent steps into shared vertices and scores a chain by
  the centrality of the vertices it visits; \textbf{gBD} forms the band from random walks on the
  reasoning graph. vBD and gBD are order-invariant.}
  \label{fig:methods_overview}
\end{figure}

\subsection{Experiments}
\paragraph{Setup.} For each question we sample $M$ Chain-of-Thought responses from a frozen LLM
with temperature sampling, parse each response into reasoning steps, embed the steps with a small
sentence encoder (all-MiniLM-L6-v2), and score every response by its band depth within the
ensemble. Confidence is evaluated as failure prediction --- the AUROC of the confidence score
against answer correctness --- on five reasoning benchmarks (GSM8K, SVAMP, ASDiv, HotpotQA,
2WikiMHQA) with Llama-3.1-8B, Llama-2-13B, and Qwen2.5-3B. We compare against three black-box baselines:
verbalized self-probing~\cite{xiong2024llms}, semantic
entropy~\cite{kuhn2023semantic,farquhar2024detecting}, and an NLI entailment score.

\paragraph{Results.} Table~\ref{tab:llm_uq} reports AUROC for Llama-3.1-8B, comparing the graph
band-depth variant (gBD) against the baselines. The reasoning-structure signal outperforms both
verbalized self-probing and answer-level consistency on every dataset: gBD raises AUROC by
$0.08$--$0.15$ over semantic entropy, with the largest gains on multi-hop QA (HotpotQA, 2WikiMHQA),
where reasoning structure --- not just the final answer --- carries most of the uncertainty. The
ordering band depth $>$ NLI / semantic entropy $>$ self-probing holds across all three models. Among our
variants gBD is best on four of five datasets, confirming that sampling reference paths from the
reasoning graph gives a richer depth estimate than the raw ensemble; the vertex variant vBD leads
on 2WikiMHQA ($0.692$). The whole family is model-free, adding no LLM calls beyond the sampled
ensemble.

\begin{table}[ht]
\centering
\caption{Failure-prediction AUROC ($\uparrow$) on five reasoning benchmarks (Llama-3.1-8B).
Baselines: verbalized self-probing, semantic entropy, and an NLI entailment score. \textbf{Bold} =
best. gBD is our graph band-depth method.}
\label{tab:llm_uq}
\small
\begin{tabular}{lcccc}
\toprule
Dataset & Self-Probing & Semantic Entropy & NLI & gBD (ours) \\
\midrule
GSM8K     & 0.528 & 0.608 & 0.637 & \textbf{0.687} \\
SVAMP     & 0.538 & 0.628 & 0.692 & \textbf{0.748} \\
ASDiv     & 0.498 & 0.671 & 0.724 & \textbf{0.761} \\
HotpotQA  & 0.549 & 0.623 & 0.638 & \textbf{0.772} \\
2WikiMHQA & 0.567 & 0.530 & 0.583 & \textbf{0.684} \\
\bottomrule
\end{tabular}
\end{table}

\subsection{Relationship to Work 1}
Both works are organized around diversity. Work~1 \emph{induces} diversity among ensemble members
during training and decides, from training dynamics, when it has been established. Work~2
\emph{reads} diversity out of an ensemble of reasoning chains at inference and turns its structure
into an uncertainty estimate. STOP's efficiency discipline carries over as a design constraint ---
Work~2 deliberately avoids the per-member training and serving cost of a true LLM ensemble by
sampling one model and staying model-free --- but the scientific object is the diversity signal
itself, not a training-dynamics trigger; there is no edge-of-stability mechanism in Work~2.

% =====================================================================
\section{Plan Towards Completion}
\label{sec:plan}

\subsection{Work~3 --- planned contribution}
With Works~1 and~2 complete, the remaining research contribution is Work~3, which continues the
LLM-uncertainty line of Work~2. \emph{Working title: adaptive, calibrated uncertainty for LLM
reasoning.} A traditional deep ensemble --- training and serving $M$ independent LLMs --- is
infeasible at LLM scale, so this line of work pursues an \emph{alternative} to deep ensembling:
approximate the ensemble by sampling multiple reasoning chains from a single model and reading
their diversity. Work~2 established that this structural diversity is a strong \emph{ranking} signal
computed from a \emph{fixed} ensemble of $M$ chains; Work~3 targets the two gaps that remain for
practical use:
\begin{itemize}[itemsep=2pt]
  \item \textbf{Adaptive sampling.} Decide \emph{how many} reasoning chains to draw per question,
  stopping once the diversity signal has converged --- an inference-time echo of STOP's
  dynamics-triggered stopping --- so easy questions cost few samples and hard ones get more.
  \item \textbf{Calibration.} Map the diversity-depth score to a \emph{calibrated} probability of
  correctness, not only a discriminative ranking, while keeping the black-box, model-free setting.
\end{itemize}
This unifies the thesis threads at the LLM level: diversity as the signal, efficiency through
adaptive sampling, and uncertainty that is both discriminative and calibrated.
\TODO{fix the precise research question, method, and evaluation once the work begins.}

\subsection{Remaining work}
\begin{enumerate}[itemsep=2pt]
  \item \textbf{Publish Works~1 and~2.} Both contributions are complete; what remains is
  write-up and publication --- STOP's revised manuscript and camera-ready, and the Work~2 paper.
  \TODO{targets}
  \item \textbf{Carry out Work~3.} Fix the research question above, implement the adaptive-sampling
  and calibration components, and run the evaluation. \TODO{key milestones}
  \item \textbf{Thesis assembly.} Integrate the three contributions --- together with the earlier
  diversity-regularization work --- into a coherent thesis with a unified introduction and a
  background chapter on ensembles, UQ, diversity regularization, and edge-of-stability dynamics.
\end{enumerate}

\subsection{Indicative timeline}
\begin{table}[ht]
  \centering
  \caption{Indicative timeline to submission. Dates to be finalized with supervisors.}
  \label{tab:timeline}
  \small
  \begin{tabular}{p{7.6cm} l}
    \toprule
    Milestone & Target \\
    \midrule
    STOP revision / camera-ready complete             & \TODO{quarter} \\
    Work~2 paper submitted                             & \TODO{quarter} \\
    Work~3 method and experiments complete             & \TODO{quarter} \\
    Work~3 paper submitted                             & \TODO{quarter} \\
    Full thesis draft to supervisors                   & \TODO{quarter} \\
    Thesis submission                                  & \TODO{quarter} \\
    \bottomrule
  \end{tabular}
\end{table}

\subsection{Risks and mitigations}
\begin{description}[leftmargin=2.4em, style=nextline, itemsep=3pt]
  \item[Scaling risk] STOP is validated at small/medium scale. \emph{Mitigation:} add at least one
  larger model / additional image domain; the monitoring cost being 2--5\% suggests the approach
  remains cheap at scale.
  \item[Regularizer-dependence risk] Benefits are less consistent for Repul. \emph{Mitigation:}
  the trajectory analysis (Fig.~\ref{fig:diversity}) is being used to characterize when STOP helps;
  report the operating regime honestly rather than overclaiming universality.
  \item[Work~3 feasibility] Work~3 is not yet started, and adaptive sampling plus calibration adds
  moving parts. \emph{Mitigation:} both components build directly on the completed Work~2 pipeline
  (model-free, no extra LLM calls), so they can be prototyped cheaply and independently; a
  fixed-ensemble, ranking-only result is already in hand from Work~2 as a fallback if calibration
  proves hard.
  \item[Publication timing] Review cycles are uncertain. \emph{Mitigation:} sequence the three
  contributions so the thesis does not depend on any single acceptance.
\end{description}

\subsection{Proposed thesis structure}
\begin{enumerate}[itemsep=1pt]
  \item Introduction --- motivation, central research question, contributions.
  \item Background --- deep ensembles, uncertainty quantification, diversity regularization, and
  the edge of stability.
  \item Diversity-aware learning --- the earlier diversity-regularization contribution
  (\emph{not covered in this progress report}). \TODO{chapter title}
  \item STOP --- stability-triggered stopping of diversity regularization (Work~1).
  \item Reasoning-diversity uncertainty for large language models (Work~2).
  \item Adaptive, calibrated uncertainty for LLM reasoning (Work~3). \TODO{confirm chapter title}
  \item Conclusion --- findings, limitations, and future directions.
\end{enumerate}

% =====================================================================
\section{Research Data Management}
\label{sec:rdm}

This section records the research data management plan (RDMP) for the project and the updates made
to it since CA1, in line with the university's research data management policy.
\TODO{cite/link the institutional RDMP record or tool entry, e.g. the Stash / DMP record}

\paragraph{Data description.}
The project uses no human-subject, personal, or otherwise sensitive data. Three kinds of data are
involved: (i) \emph{public benchmark datasets} --- UCI regression sets and CIFAR-10 for Work~1, and
GSM8K, SVAMP, ASDiv, HotpotQA, and 2WikiMHQA for Work~2 --- all obtained from their original public
releases under their respective licenses; (ii) \emph{generated data} --- ensemble members and
training logs (Work~1), and sampled Chain-of-Thought responses, step embeddings, per-question
confidence scores, and evaluation metrics (Work~2); and (iii) \emph{code and configuration} that
produce the above. Pre-trained language models (Llama~2 / Llama~3 / Qwen2.5) are used under their community
licenses and are treated as inputs, not redistributed.

\paragraph{Storage, backup, and organization.}
All source code is under Git version control, with each contribution in its own repository and a
documented directory layout (raw inputs, generated outputs, and analysis kept separate). Working
data and model checkpoints are held on institutional research storage with automatic backup, rather
than on local machines. \TODO{name the storage location / quota and the backup cadence}
Generated artifacts are organized by \texttt{model / dataset / method} so that any reported number
is traceable to the script and configuration that produced it.

\paragraph{Retention, sharing, and preservation.}
Because all inputs are public and no sensitive data is involved, the intended default is open
release: code and the scripts needed to regenerate results will be published in a public repository
on paper acceptance, and derived result files (metrics, confidence scores) archived alongside.
Large regenerable artifacts (checkpoints, raw sampled responses) are retained for the project's
duration and for the university's required minimum retention period, then reviewed for deletion.
\TODO{state the retention period and the intended archive/DOI (e.g. Zenodo / institutional repository)}

\paragraph{Updates and active engagement since CA1.}
The RDMP has been kept current with the project rather than treated as a one-off form:
\begin{itemize}[itemsep=2pt]
  \item Work~2 introduced a new class of generated data (sampled CoT ensembles, embeddings,
  confidence scores); the plan was extended to cover its layout, storage, and provenance.
  \item Adopted a consistent, reproducible pipeline (fixed seeds, resume-safe runs, and an
  auto-derived output path per \texttt{model/dataset}) so results are regenerable from code plus
  public data.
  \item \TODO{list any concrete engagement since CA1 --- e.g. migrated working data to institutional
  storage, updated the DMP record on <date>, confirmed dataset licenses.}
\end{itemize}
No ethics approval is required, as the work uses only public datasets and pre-trained models and
generates no data about human participants.

\bibliographystyle{plain}
\bibliography{refs}

\end{document}
